\begin{resumo}[Resumo]
\begin{otherlanguage*}{brazil}
Este trabalho desenvolve e avalia um simulador modular em Python do pipeline \emph{DNA Fountain}, com foco nos códigos LT (\emph{Luby Transform}), na Distribuição Robusta de Soliton (RSD) e no algoritmo de decodificação iterativa \emph{peeling} em regime de blocos finitos ($K \le 64$). O simulador modela a segmentação de cargas binárias, a síntese de gotas, a filtragem por restrições físico-químicas de \emph{screening}, perdas probabilísticas em canal de apagamento e a decodificação em grafos bipartidos sobre o corpo finito binário $GF(2)$. Uma avaliação fatorial com 216 combinações paramétricas e 1.000 repetições independentes com sementes sorteadas (216.000 ensaios) caracteriza o desempenho sob diferentes taxas de redundância e perdas. Sob canal ideal e \emph{overhead} $\epsilon = 1{,}00$, a configuração de referência ($c = 0{,}1$; $\delta = 0{,}10$) alcançou $92{,}0\%$ de sucesso na recuperação para $K = 32$ (intervalo de confiança exato de Clopper--Pearson de 95\%: $[90{,}1\%; 93{,}6\%]$) e $98{,}5\%$ para $K = 64$ ($[97{,}5\%; 99{,}2\%]$). Para $K = 16$, a parametrização $c = 0{,}3$ superou $c = 0{,}1$ ($89{,}3\%$ contra $84{,}7\%$; $p = 0{,}0027$), apresentando maior recuperação também sob orçamentos de redundância restritos ($\epsilon \le 0{,}50$), o que indica que nenhuma parametrização isolada domina o regime de blocos finitos. Sob canal ideal e $\epsilon = 1{,}00$, a configuração $c = 0{,}5$ apresentou $38{,}5\%$ de sucesso para $K = 32$ em razão do colapso do modo corretivo para $m = \lfloor K/R \rfloor = 1$ quando $R > K/2$, enquanto para $K = 64$ ($R < K/2$) sustentou $69{,}2\%$ de resolução em virtude do deslocamento modal para $m = 2$. A auditoria de falhas em $K = 16$ com validação cruzada por eliminação gaussiana em $GF(2)$ indicou que $98{,}7\%$ dos insucessos decorreram de conjuntos de parada (\emph{stopping sets}) cíclicos com posto completo de matriz, alcançando o decodificador algébrico $99{,}8\%$ de resolução. O \emph{screening} biológico rejeitou $42{,}6\%$ dos candidatos enquadrados (referência analítica exata de $49{,}8\%$ para 28-meros uniformes não restritos), sem diferença estatisticamente significativa entre as distribuições de graus de candidatos aceitos e rejeitados (qui-quadrado de homogeneidade = 23,66; 17 graus de liberdade; $p = 0{,}1289$). Sob probabilidades de apagamento de 0,10 e 0,20 no canal com \emph{overhead} nominal $\epsilon = 1{,}00$, a recuperação decresceu para $82{,}9\%$ e $63{,}4\%$ ($K = 32$) e para $92{,}3\%$ e $78{,}8\%$ ($K = 64$), em conformidade com a redução do número efetivo de gotas recebidas. O estudo limita-se à modelagem computacional e simulação em software, sem contemplar síntese química ou validação física em bancada molecular.

\vspace{\onelineskip}

\noindent
\textbf{Palavras-chave}: Armazenamento de dados em DNA; \emph{DNA Fountain}; Códigos LT; Decodificação \emph{peeling}; Simulação de software.
\end{otherlanguage*}
\end{resumo}
